\par\Needspace{8\baselineskip}\bigskip
\SourceLeaf{487}%
\pdfbookmark[0]{Entwurf zur Gradmessung. 30. Mai 1819.}{entwurf-gradmessung-mai1819}
\begin{center}\large [ENTWURF ZUR GRADMESSUNG.]\end{center}
\par
— — — Über Gradmessungen im Allgemeinen und darüber, inwiefern sie einen der interessantesten Gegenstände menschlichen Wissens bilden, brauche ich in einem Schreiben an Ew. — — nichts zu sagen. Die großen Vorteile und die neue Aufklärung noch dunkler Punkte, die von der Vervielfachung solcher Unternehmungen zu erwarten sind, hängen aber auch davon ab, dass diese in möglichst großem Maßstab durchgeführt werden. Vereinzelte Gradmessungen in Europa, die nur einen kleinen Bereich umfassen, könnten jetzt, nach den großen Arbeiten in Frankreich und England, nur einen sehr untergeordneten Wert haben. Dagegen wären noch eine oder ein paar weitere ähnliche Messungen von beträchtlicher Ausdehnung in Europa für die Kenntnis der Gestalt der Erde gewiss ungemein wichtig.

Ein Blick auf die Karte Europas und auf den Kulturstand der verschiedenen Nationen zeigt, dass außer der großen Linie von den Balearen bis zu den Orkneyinseln nur noch in zwei Richtungen ähnliche Unternehmungen durchführbar sein werden, 1) im Russischen Reich und 2) durch die Halbinsel Jütland und ganz Deutschland bis zum Mittelmeer. Nach den Nachrichten, die mir zugegangen sind, scheint es einige Aussichten auf die Ausführung einer Messung im erstgenannten Reich zu geben. Für die andere gibt es aber mehr als eine Aussicht. Der erste wesentliche Schritt ist bereits tatsächlich geschehen. Die vom König von Dänemark angeordnete Gradmessung von der Nordspitze Jütlands bis Lauenburg ist bereits seit zwei Jahren im Gange. Dafür, dass sie ganz so ausgeführt wird, wie es der heutige Stand der Wissenschaften und Künste möglich und notwendig macht, bürgen die Geschicklichkeit und Einsicht des dänischen Astronomen, die Vortrefflichkeit seiner Instrumente und die Großzügigkeit, mit der der König von Dänemark alles genehmigt, was zur Vollkommenheit der Messung für notwendig oder wünschenswert gehalten wird.

Die dänische Gradmessung soll außer dem erwähnten Meridianbogen auch die Messung eines Breitenkreisbogens von der Westküste Jütlands bis Kopenhagen umfassen. Natürlich kann hier nur vom ersten Teil des Plans die Rede sein. Jener Bogen umfasst für sich genommen schon \(4\frac12\) Grad, und die Messung ist daher auch für sich betrachtet von beträchtlicher Ausdehnung. Ungleich wichtiger erscheint sie aber, wenn man sie als Anfang jener großen Messung betrachtet, die sich bis zur Insel Elba, also bis auf \(16^0\), ausdehnen ließe. Dass diese Unternehmung früher oder später einmal in einer solchen Ausdehnung ausgeführt wird, ist wohl mehr als eine wirklichkeitsferne Hoffnung, zumal in einigen dazwischenliegenden Ländern, namentlich in Thüringen und Bayern, tatsächlich schon manche Vorarbeiten geleistet worden sind. Die Fortsetzung der dänischen Messung durch das Königreich Hannover ist aber die erste und wesentlichste Voraussetzung für die spätere Verwirklichung jenes Plans. Schon durch diese Fortsetzung allein würde der Bogen um 2 Grad vergrößert. Würde dann auch noch die Gothaer Sternwarte durch ein Dreiecksnetz mit der Göttinger verbunden, was auch in mancher anderen Hinsicht sehr wünschenswert und leicht durchführbar wäre, so wäre damit schon ein Bogen von 7 Grad verwirklicht.

Ich brauche nur kurz zu erwähnen, dass die Messung eines Meridianbogens von Hamburg bis Göttingen auch in anderer als rein wissenschaftlicher Hinsicht von großer Bedeutung wäre. Falls über kurz oder lang eine den heutigen Anforderungen entsprechende Vermessung des ganzen Königreichs Hannover beschlossen werden sollte, würde das zu diesem Zweck angelegte Dreiecksnetz die sicherste Grundlage bieten, um die weiteren Triangulationen nach Osten und Westen daran anzuschließen. Und falls eine solche Gesamtvermessung in naher Aussicht stehen sollte, ließe sich mit der Gradmessung noch der Nebenzweck erreichen, dass sie auch zur Vorbereitung geeigneter Personen auf jene Arbeit dienen könnte.
